% from ega4-9-fr.tex
\oldpage[IV-3]{92}

composantes irréductibles de $Y_\eta$, il résulte de (9.7.7), (9.7.8) et
(9.5.3) que pour tout $s$ voisin de $\eta$, les $(Z_i)_s$ sont les composantes
irréductibles de $Y_s$ et qu'elles sont \emph{géométriquement irréductibles}.
Dire qu'en un point $y\in X_s$, $\mathcal F_s$ est géométriquement intègre
signifie alors que $\mathcal F_s$ est géométriquement réduit en ce point et en
outre que $y$ n'appartient qu'à \emph{un seul} des $(Z_i)_s$ (4.6.22). La
conclusion résulte donc d'une part de (iv) et d'autre part de (9.5.1) appliqué à
l'intersection de $\overline{\{x\}}$ et de chaque $Z_i$.

(vi) Gardant les mêmes notations que dans (ii), il résulte de (6.2.1) que l'on
a
\[
\operatorname{dim.proj}((\mathcal F_s)_x)=
\operatorname{dim.proj}(((\mathcal F_0)_{s_0})_{x_0})~;
\]
on peut donc encore se borner au cas où $A$ est \emph{noethérien}. En outre, on
se ramène encore à montrer que, pour tout $x\in X$, il existe un voisinage $V$
de $x$ dans $\overline{\{x\}}$ tel que, pour tout $x'\in V$, on ait
\[
\operatorname{dim.proj}((\mathcal F_{f(x')})_{x'})=
\operatorname{dim.proj}((\mathcal F_{f(x)})_x)~;
\]
et comme ci-dessus, on peut supposer que $S$ est \emph{intègre} et que
$\eta=f(x)$ est son point générique, de sorte que l'on a
$(\mathcal F_\eta)_x=\mathcal F_x$. En vertu du théorème de platitude générique
(6.9.1), on peut, en remplaçant au besoin $S$ par un voisinage ouvert de
$\eta$, supposer que le morphisme $f$ est \emph{plat} et que $\mathcal F$ est
$f$-plat~; on a alors
$\operatorname{dim.proj}((\mathcal F_{f(x')})_{x'})=
\operatorname{dim.proj}(\mathcal F_{x'})$ pour tout $x'\in X$ en vertu de
(6.2.3). Cela étant, en vertu de (6.11.1), on peut (en remplaçant au besoin $X$
par un voisinage ouvert de $x$) supposer que
$\operatorname{dim.proj}(\mathcal F_{x'})\leq
\operatorname{dim.proj}(\mathcal F_x)$ pour tout $x'\in X$. D'autre part, si
$\operatorname{dim.proj}(\mathcal F_x)=n$, il y a par hypothèse un
$\mathcal O_x$-module de type fini $M$ tel que
$\operatorname{Ext}_{\mathcal O_x}^n(\mathcal F_x,M)\neq0$
($\mathbf 0$, 17.2.4). Or, il existe un $\mathcal O_X$-Module cohérent
$\mathcal G$ tel que $M=\mathcal G_x$ (en remplaçant au besoin $X$ par un
voisinage ouvert de $x$ ($\mathbf 0_{\mathrm I}$, 5.3.8))~; en vertu de
(T, 4.2.2), on a donc
$(\mathcal{E}xt_{\mathcal O_X}^n(\mathcal F,\mathcal G))_x\neq0$. Mais
$\mathcal{E}xt_{\mathcal O_X}^n(\mathcal F,\mathcal G)$ est un
$\mathcal O_X$-Module cohérent ($\mathbf 0_{\mathrm{III}}$, 12.3.3), donc son
support $Y$ est fermé ($\mathbf 0_{\mathrm I}$, 5.2.2)~; comme il contient $x$,
il contient aussi $\overline{\{x\}}$, d'où l'on conclut (en appliquant encore
(T, 4.2.2)) que l'on a
$\operatorname{dim.proj}(\mathcal F_{x'})\geq n$ pour tout
$x'\in\overline{\{x\}}$, ce qui achève de prouver l'assertion dans le cas (vi).

(vii) Comme $B$ est une $A$-algèbre de type fini, $X$ est $S$-isomorphe à un
\emph{sous-schéma fermé} d'un $S$-schéma de la forme
$Y=\operatorname{Spec}(A[T_1,\ldots,T_r])$~; si $j:X\to Y$ est l'injection
canonique, et $\mathcal G=j_*(\mathcal F)$, on a
$\mathcal G_s=(j_s)_*(\mathcal F_s)$ pour tout $s\in S$, et, en vertu de
($\mathbf 0$, 16.4.11), on peut se borner à démontrer l'assertion pour $Y$ et
$\mathcal G$. Autrement dit, on peut supposer que $B=A[T_1,\ldots,T_r]$, de
sorte que chacun des schémas $X_s=\operatorname{Spec}(k(s)[T_1,\ldots,T_r])$
est \emph{régulier} ($\mathbf 0$, 17.3.7). Soit alors
$W=\operatorname{Supp}(\mathcal F)$, de sorte que
$W_s=\operatorname{Supp}(\mathcal F_s)$ (\textbf{I}, 9.1.13)~; on a, par
(6.11.2.1)
\begin{equation}
\tag{9.9.2.1}
\label{IV.9.9.2.1-fr}
\operatorname{coprof}((\mathcal F_{f(x)})_x)=
\operatorname{dim.proj}((\mathcal F_{f(x)})_x)-
\operatorname{codim}_x(W_{f(x)},X_{f(x)}).
\end{equation}
Mais comme $W$ est constructible (8.9.1) et chaque $W_s$ fermé, il résulte de
(vi) et de (9.9.1, (ii)) que les deux fonctions du second membre de (9.9.2.1)
sont constructibles~; il en est donc de même de leur différence ce qui achève
de prouver la proposition dans le cas (vii).

(viii) Soit $U_n$ l'ensemble des $x\in X$ tels que
$\operatorname{coprof}((\mathcal F_{f(x)})_x)\leq n$, et posons
$Z_n=X-U_n$~; il résulte de (vii) que les $Z_n$ sont constructibles~; en outre,
comme

